\documentclass[12pt]{article}
\usepackage{fontspec}
\usepackage{polyglossia}
\usepackage{amsmath}
\usepackage{amssymb}
\usepackage[margin=2.5cm]{geometry}
\usepackage{tikz}
\usepackage{luacode}

\setmainlanguage{turkish}
\setotherlanguage{english}

\setmainfont{Arial}

\title{LuaLaTeX Yetenek Testi}
\author{Sabir Aliyev}
\date{\today}

\begin{document}

\maketitle

\section{Unicode ve Coklu Dil Destegi}
Turkce ozel karakterler: ğ, ş, ı, ö, ü, ç, İ, Ğ, Ş, Ö, Ü, Ç dogrudan yaziliyor,
ekstra paket (inputenc/babel) gerekmiyor.

\begin{english}
English text mixed in: The quick brown fox jumps over the lazy dog.
\end{english}

Gurcuce (Unicode, farkli bir alfabe):
\begin{center}
{\fontspec{Noto Sans Georgian} გამარჯობა მსოფლიო, ეს არის ტესტი}
\end{center}

\section{Lua Kodu Calistirma}
LuaLaTeX'in en guclu ozelligi: belge icinde gercek Lua kodu calistirabilmek.
Asagidaki hesaplama Lua ile yapildi, LaTeX degil:

\begin{luacode}
function fibonacci(n)
  local a, b = 0, 1
  for i = 1, n do
    a, b = b, a + b
  end
  return a
end

tex.print("10. Fibonacci sayisi: " .. fibonacci(10))
\end{luacode}

\section{Matematik}
\begin{equation}
    \sum_{n=1}^{\infty} \frac{1}{n^2} = \frac{\pi^2}{6}
\end{equation}

\begin{equation}
    e^{i\pi} + 1 = 0
\end{equation}

\section{Tikz Grafik}
\begin{center}
\begin{tikzpicture}
\draw[thick, blue] (0,0) circle (2cm);
\draw[thick, red] (-2,0) -- (2,0);
\draw[thick, red] (0,-2) -- (0,2);
\node at (0,2.5) {LuaLaTeX + Tikz};
\end{tikzpicture}
\end{center}

\section{Sonuc}
Eger bu belge sorunsuz derlendiyse, backend artik:
\begin{itemize}
    \item Unicode karakterleri (Turkce, Gurcuce) ekstra ayar gerekmeden,
    \item Coklu dil desteğini (\texttt{polyglossia}),
    \item Belge icinde calisan gercek Lua kodunu (\texttt{luacode}),
\end{itemize}
destekliyor demektir.

\end{document}